\section{First-order sequent rules}
\label{sec:fol-deduction-02}

The first-order sequent calculus \(\LK_{\FOL}\) extends
\(\LK_{\CPL}\).  A sequent is satisfied by \(\Meta,s\) under the same
left-conjunctive and right-disjunctive reading used for propositional
sequents.  Its additional rules are
\[
\begin{array}{c@{\qquad}c}
\Rule[\(\forall L\)]{\Gamma,\varphi[t/x]\Rightarrow\Delta}
 {\Gamma,\forall x\,\varphi\Rightarrow\Delta}
&
\Rule[\(\exists R\)]{\Gamma\Rightarrow\Delta,\varphi[t/x]}
 {\Gamma\Rightarrow\Delta,\exists x\,\varphi}
\\[2.8ex]
\Rule[\(\forall R\)]{\Gamma\Rightarrow\Delta,\varphi[y/x]}
 {\Gamma\Rightarrow\Delta,\forall x\,\varphi}
&
\Rule[\(\exists L\)]{\Gamma,\varphi[y/x]\Rightarrow\Delta}
 {\Gamma,\exists x\,\varphi\Rightarrow\Delta}.
\end{array}
\]
In \(\forall R\), the eigenvariable \(y\) occurs free in no formula of
\(\Gamma,\Delta,\forall x\,\varphi\).  In \(\exists L\), it occurs free in
no formula of \(\Gamma,\Delta,\exists x\,\varphi\).

Equality may be included by the initial sequent
\(\Rightarrow t=t\) and the replacement rule
\[
 \Rule[\(=L\)]{
   \Gamma,\varphi[r/x]\Rightarrow\Delta}
  {\Gamma,r=t,\varphi[t/x]\Rightarrow\Delta}.
\]
The symmetric replacement instance---from
\(\Gamma,\varphi[t/x]\Rightarrow\Delta\) infer
\(\Gamma,r=t,\varphi[r/x]\Rightarrow\Delta\)---is derivable.  First,
\(r=t\Rightarrow t=r\) is derivable: apply \(=L\) with the auxiliary
formula \(t=x\), for a variable \(x\) absent from \(r\) and \(t\), to
the identity sequent \(t=r\Rightarrow t=r\), obtaining
\(r=t,\,t=t\Rightarrow t=r\), and cut with the initial sequent
\(\Rightarrow t=t\).  Second, the displayed rule with the roles of the two
terms exchanged turns the premise into
\(\Gamma,t=r,\varphi[r/x]\Rightarrow\Delta\), and a cut with
\(r=t\Rightarrow t=r\) yields the instance.  A first-order sequent
derivation is a
finite tree generated by the propositional and displayed first-order rules.

